% SPDX-FileCopyrightText: 2026 Will Cook
% SPDX-License-Identifier: CC-BY-4.0
%
% Standalone Erdős Problem Note for #68.  Context is deliberately repeated
% from the sibling notes so that this file remains independently readable.
%
% Build from the public paper directory with `make`.
\documentclass[11pt]{article}
\input{problem-note-preamble}
\usepackage{longtable}
\hypersetup{
  pdftitle={Two incomparable denominator exclusions for the sum of reciprocals of n!-1 (Erdős Problem 68)},
  pdfsubject={Erdős Problem Notes: Problem 68},
  pdfkeywords={irrationality, factorial denominators, companion constant, factorial orbit, strict successor, prime-pole residue, continued fractions, Lean 4},
}

% The prose keeps compact declaration labels; the source appendix supplies the
% commit-pinned links checked by the public release gate.
\newcommand{\pdecl}[1]{\emph{\leanlabel{#1}}}
\newcommand{\Ser}{S}
\newcommand{\chn}[2]{V_{#1}(#2)}
\newcommand{\Lc}{L}

\runninghead{Erd\H{o}s \#68: two incomparable denominator exclusions}
\title{Two Incomparable Denominator Exclusions\\for $\sum_{n\ge2}(n!-1)^{-1}$}
\subtitle{Finite certificates and exact rationality criteria for Erd\H{o}s Problem~\#68}
\author{Will Cook}
\date{31 August 2026}

\begin{document}
\maketitle
\noteprovenance

\begin{abstract}
\noindent
Every displayed rational representation $\Ser=a/q$ with $q>0$, where
\[
  \Ser=\sum_{n\ge2}\frac1{n!-1},
\]
satisfies two incomparable exclusions.  The first is a divisibility exclusion,
\[
  q\nmid299999! ,
\]
so some prime power in $q$ exceeds its multiplicity in $299999!$.  The
second is a size exclusion,
\[
  q\ge2^{39990}>10^{12038}.
\]
Together these restrict the denominator's size and prime-power multiplicities.
Neither restriction implies the other.  The first rests on a Lean-checked consumer together with one exact
GMP carry computation, $300000\nmid Z_{300000}$; the second on $23449$
certified continued-fraction partial quotients from an $80000$-bit exact
rational bracket.  Both computations use exact integer arithmetic.

Two consecutive unit carries reduce, after an exact denominator telescope, to
the universal inequality $0<U_{m+2}\le V_{m+2}$.  That cleared window adds no
contradiction after the two carries have been imposed; it does not quantify
over every possible predecessor restriction.

The exclusions are read off exact coordinates for the problem, which we state
next because they fix the missing quantifier rather than because they advance
it.  Put $H_m=\sum_{2\le n\le m}(n!-1)^{-1}$ and $Z_m=\lfloor m!H_m\rfloor+1$,
and let $b_m$ be the exact integer carry in the recurrence
$Z_m=mZ_{m-1}+1-b_m$.  Lean proves that $\Ser$ is rational precisely when
$b_m=1$ for all sufficiently large $m$, and $b_m=1$ is equivalent to $m\mid
Z_m$ for every $m\ge3$, so
\[
  \Ser\notin\mathbb Q
  \quad\Longleftrightarrow\quad
  (\forall B)(\exists m>B)\;m\nmid Z_m .
\]
A companion series gives a second coordinate.  With
$C=\sum_{n\ge2}(n!(n!-1))^{-1}$ the termwise identity $C+(e-2)=\Ser$ yields
\[
  \Ser\in\mathbb Q
  \quad\Longleftrightarrow\quad
  \lfloor m!C\rfloor\equiv-2\pmod m
  \quad\hbox{for all sufficiently large }m,
\]
so irrationality of $\Ser$ is cofinal escape of one explicit factorial orbit
from one distinguished residue class.  Neither equivalence proves that the
escape occurs.

A further reduction describes cancellation at a fixed prime.  If $q^e$
is the largest power of a prime $q$ occurring among the
displayed denominators $n!-1$, then the common-denominator numerator modulo $q$ is
a $q$-adic unit times the sum of the reciprocal maximal-hit cofactors.  Hence
the complete power $q^e$ survives prefix reduction exactly when this
principal residue is nonzero.  Complete cancellation occurs at the Wilson
prefixes for $q=139$ and $q=2593$; the corresponding residue identities are
also checked in Lean.  The set of maximal hits therefore does not determine
whether $q^e$ survives reduction: their cofactor residues matter.

The remaining problem is to prove $m\nmid Z_m$ for arbitrarily large $m$.
None of the finite computations, exclusions, or structural reductions in this
paper proves that assertion.
\end{abstract}

\section{The problem}\label{sec:problem}

\begin{problem}[Erd\H{o}s \#68]\label{res:problem}
Is
\[
  \Ser=\sum_{n\ge2}\frac1{n!-1}
\]
irrational?
\end{problem}

\subsection*{Main results at a glance}

\begin{center}
\small
\begin{tabular}{@{}
  >{\raggedright\arraybackslash}p{0.23\linewidth}
  >{\raggedright\arraybackslash}p{0.69\linewidth}@{}}
\toprule
Result & Exact scope \\
\midrule
\textbf{Fixed companion-orbit boundary} & For $C=\sum_{n\ge2}(n!(n!-1))^{-1}$, one has $\Ser\in\mathbb Q$ iff $\lfloor m!C\rfloor\equiv-2\pmod m$ eventually; equivalently, irrationality is cofinal escape from that residue. \\
\textbf{Exact irrationality criterion} & $\Ser\notin\mathbb Q$ iff for every $B$ some $m>B$ satisfies $m\nmid Z_m$. \\
\textbf{Finite denominator exclusions} & Two incomparable statements: $q\nmid299999!$ from the carry certificate and $q\ge2^{39990}>10^{12038}$ from certified continued fractions. \\
\textbf{Prime-pole cancellation law} & For a finite prefix and prime $q$ with positive attained maximal exponent $e$, the complete power $q^e$ survives reduction iff the reciprocal sum of its maximal-hit cofactors is nonzero.  Exact scans give full cancellations at $q=139$ and $2593$; Lean checks the residue equalities. \\
\textbf{Adjacent-unit obstruction} & After the pair assumption, the complete width-one window cancels to the universal reduced-numerator bound; it cannot exclude the pair. \\
\textbf{Elementary lcm growth} & $\liminf\log L_N/(N^{3/2}\log N)\ge 2\sqrt2/3$ is Stirling packaging of a Lean-checked terminal-block inequality, not a reduced-denominator bound.  The same coefficient holds for $\operatorname{lcm}(n!+P(n))$ when $P\in\mathbb Z[X]$ is fixed and nonzero; $P=0$ collapses to $N!$. \\
\textbf{Moving-factor criterion} & One moving prefix-private prime plus two explicit scale bounds implies irrationality; the first bound is already the registered global complementary-residue consumer.  Arbitrary split factors give a normalized-collision criterion at scale $2p^2(2p-1)!/(p-1)!$.  Every fixed owner pair is eventually absorbed.  The cofinal scale and disagreement producers remain open. \\
\textbf{Further reductions} & Checked channel identities and projection lemmas state the additional estimates that those methods would require. \\
\textbf{Open step} & No cofinal miss or cofinal strict residual nonvanishing is proved; the irrationality question remains open. \\
\bottomrule
\end{tabular}
\end{center}

\noindent Erd\H{o}s states the problem on p.~102 of his 1988 survey and, in
the same passage, records the broader expectation that
$\sum_n1/(n!+t)$ is irrational, indeed transcendental, for every integer
$t$ \cite[p.~102]{erdos1988}.  Erd\H{o}s presents this broader statement as a
conjecture.

\noindent Numbering and current status follow
\href{https://www.erdosproblems.com/68}{Bloom's Erd\H{o}s problem catalogue}
\cite{bloom}.
The problem is open.  The companion series $\sum_{n\ge0}1/n!=e$ and
$\sum_{n\ge2}1/(n!+1)$ sit in the same family, and the difficulty here is the
same one that makes the Erd\H{o}s--Borwein constant hard: the denominators
$n!-1$ grow fast enough that convergence is trivial and slow enough, in the
arithmetic sense, that no single congruence controls them.

This note is one of the problem-owned series.  It repeats the definitions and
the claim boundary it needs, so that a reader who arrives at \#68 alone does
not have to assemble them from the joint integration monograph.

\statusboundary

\small
\begin{longtable}{@{}
  >{\raggedright\arraybackslash}p{0.34\linewidth}
  >{\raggedright\arraybackslash}p{0.20\linewidth}
  >{\raggedright\arraybackslash}p{0.38\linewidth}@{}}
\toprule
Statement & Status & Exact boundary \\
\midrule
\endfirsthead
\toprule
Statement & Status & Exact boundary \\
\midrule
\endhead
\bottomrule
\endfoot
Irrationality of $\Ser$ & Open & No proof is claimed. \\
Fixed companion-orbit boundary & Checked & The two endpoint equivalences are checked in \path{CompanionOrbitRationality.lean}; their exact declaration names and source routes appear in Appendix~\ref{app:sources}.  Both depend exactly on \texttt{propext}, \texttt{Classical.choice}, and \texttt{Quot.sound}; neither proves that the cofinal misses occur. \\
Canonical factorial digit kernel & Checked & Floor formula, digit bounds, remainder recurrence, finite expansion, zero-tail propagation. \\
Channel integrality & Checked & $(d!)^{\lfloor i/d\rfloor}\mid i!$, with exact denominator cancellation. \\
Channel congruence and LCM tax & Checked & $\chn d\lambda\equiv M\pmod{d!-1}$; annihilating channels through $D$ forces $\Lc_D\mid M$. \\
Prime unit translator & Checked & The pair $(p,-1)$ on $(p-1,p)$ has moment $0$, all channels $0$ except $d=p$, and $p$-channel numerator $p!-1$. \\
Weighted projection rigidity & Checked & If $Z\equiv T\pmod R$, $Q_i\mid R$, and $Z\le B<Q_i$, then unequal residues $T\bmod Q_1$ and $T\bmod Q_2$ exclude that endpoint. \\
Factor-split projection reduction & Checked & Two divisor factors of one private modulus support the same cancellation and disagreement bounds; coprime factors give a branch-free floor and may lie in one private quotient. \\
Zero plateau and first-exit carry & Checked & Grid threshold, plateau equality of grid integers, forced zero digit, carry $b\in\{0,-1\}$. \\
Prime-power prefix obstruction & Checked & Rationality forces $p^k$ to divide the strict successor at $kp$ whenever the factorial clears the denominator and $p$ is coprime to it. \\
Exact carry characterization & Checked & The normalized strict successors converge to $\Ser$; $\Ser$ is rational exactly when $b_m=1$ eventually, equivalently $\Ser$ is irrational exactly when non-unit carries occur cofinally. \\
Explicit denominator exclusions & Checked implications; exact finite certificates & Exact reduction gives $60\nmid Z_{60}$, $64\nmid Z_{64}$, and $67\nmid Z_{67}$.  The GMP carry computation gives $q\nmid299999!$ (hence $q\ge300000$); an independent bracket-certified continued-fraction computation gives $q\ge2^{39990}>10^{12038}$.  Neither changes a quantifier. \\
Adjacent-unit cancellation normalizer & Checked & Under consecutive unit carries, the exact offset and denominator share the same positive two-step normalizer, so their window inequality is equivalent to the universal bound $0<U_{m+2}\le V_{m+2}$. \\
Digits eventually zero $\iff$ $\Ser$ rational & Checked & \path{CanonicalFactorialTermination.lean} proves the equivalence for every real; it does not decide whether $\Ser$ is rational. \\
Elementary lcm growth $\liminf\log L_N/(N^{3/2}\log N)\ge 2\sqrt2/3$ & Ordinary packaging of a Lean-checked segment inequality & Finite product/lcm/pairwise-gcd bound is kernel-checked; the Stirling liminf is an ordinary corollary.  It is not a reduced-denominator bound and is not an irrationality theorem. \\
Polynomial-shift lcm growth & Ordinary proof; gcd source landed & Same leading coefficient for $\operatorname{lcm}(n!+P(n))$, $P\in\mathbb Z[X]$ fixed and nonzero.  $P=0$ collapses to $N!$.  The gcd identity is source-landed, not a Lean receipt in this wave. \\
GLS multiplicity $\log L_N\gg N^{4/3}\log N$ & Derived, source-verified & Weaker exponent from Garaev--Luca--Shparlinski Theorem~12; retained as a cited alternative. \\
Finite certificates ($\lambda=2e_3-e_4$, $D=3$ on $(3,4,5,6)$) & Verified finite instances & Fully specified arrays only. The printed $D\le12$ moment-$L_D$ family is withdrawn on support $n\ge2$. \\
Unbounded strict nonvanishing & Open & Required to turn the channel rounding argument into an irrationality proof. \\
\end{longtable}

\section{The fixed companion orbit}\label{sec:companion-orbit}

The exact reduction of this section is a transport of a classical criterion,
and the classical criterion should be named first.  A real number is rational
exactly when its canonical factorial-base digits vanish from some index on;
this is due to Cantor~\cite{cantor1869}, and the standard modern account is
Galambos~\cite[Ch.~1]{galambos1976}.  One direction is the observation that
$S=a/b$ makes $m!S$ an integer for $m\ge b$, and the other is a tail bound.
What is added below is the transport of that criterion along an explicit
companion telescope, the identification of the resulting digit value, and the
observation that the same shape survives for the whole family
$\sum_n(n!+t)^{-1}$ after a sign adjustment.  The difficulty is entirely in the
unproved cofinal-miss producer, so this is a reduction and not an advance on
the endpoint.

The reduction uses one fixed real number rather
than the changing rational prefixes.  Define
\[
 C=\sum_{n\ge2}\frac1{n!(n!-1)}.
\]
The elementary identity
\[
 \frac1{n!-1}=\frac1{n!}+\frac1{n!(n!-1)}
\]
and absolute convergence give
\begin{equation}
 C+(e-2)=\Ser.                                      \label{eq:companion-decomposition}
\end{equation}

\begin{theorem}[fixed companion-orbit rationality boundary]
\label{res:companion-orbit-rationality-boundary}
The following statements are equivalent:
\begin{enumerate}
\item $\Ser\in\mathbb Q$;
\item $(\lfloor m!C\rfloor+2)\bmod m=0$ for every sufficiently large $m$.
\end{enumerate}
Consequently,
\begin{equation}
 \Ser\notin\mathbb Q
 \quad\Longleftrightarrow\quad
 (\forall B)(\exists m>B)\;
   (\lfloor m!C\rfloor+2)\bmod m\ne0 .
 \label{eq:companion-cofinal}
\end{equation}
\end{theorem}

\begin{proof}
First suppose $\Ser=a/q$ with $q>0$, and take $m$ so large that
$q\mid(m-1)!$.  Write
\[
 E_m=m!\sum_{2\le n\le m}\frac1{n!},
 \qquad
 \varepsilon_m=m!\sum_{n>m}\frac1{n!}.
\]
Then $E_m$ is an integer, $0<\varepsilon_m<1$, and $m!\Ser$ is an integer
divisible by $m$.  Multiplying \eqref{eq:companion-decomposition} by $m!$
therefore gives
\[
 \lfloor m!C\rfloor=m!\Ser-E_m-1.
\]
Every summand of $E_m$ except the endpoint $m!/m!=1$ is divisible by $m$.
Thus $E_m\equiv1\pmod m$ and
$\lfloor m!C\rfloor\equiv-2\pmod m$.

Conversely, assume the displayed congruence from some index onward.  Let
$d_m(C)=\lfloor m!C\rfloor-m\lfloor(m-1)!C\rfloor$ be the canonical
factorial digit.  Since $0\le d_m(C)<m$, for $m\ge3$ the congruence is
equivalent to
\[
 d_m(C)=m-2.
\]
Choose $N$ beyond the exceptional indices.  The canonical factorial
expansion of $C$ then has the form
\[
 C=\lfloor C\rfloor+
   \sum_{m=2}^{N}\frac{d_m(C)}{m!}+
   \sum_{m>N}\frac{m-2}{m!}.
\]
Adding $e-2=\sum_{m\ge2}1/m!$ and using the telescoping identity
\[
 \sum_{m>N}\frac{m-1}{m!}
 =\sum_{m>N}\left(\frac1{(m-1)!}-\frac1{m!}\right)
 =\frac1{N!}
\]
gives
\[
 \Ser=\lfloor C\rfloor+
   \sum_{m=2}^{N}\frac{d_m(C)+1}{m!}+\frac1{N!},
\]
which is rational.  Negating the eventual statement yields the cofinal
formulation above.
\end{proof}

\begin{remark}[the exact remaining boundary]
\label{bdry:companion-orbit-nonconcentration}
The theorem isolates the problem without solving it.  A proof of Problem~\#68
must still show that the orbit $\lfloor m!C\rfloor\bmod m$ does not eventually
concentrate at $-2$.  Finite computation can measure this concentration but
cannot discharge the cofinal quantifier.
\end{remark}

\section{Canonical factorial digits}\label{sec:digits}

Write $\theta_0=\{x\}=x-\lfloor x\rfloor$ and, for $m\ge1$,
\[
  d_m=\lfloor m\,\theta_{m-1}\rfloor,
  \qquad
  \theta_m=m\,\theta_{m-1}-d_m .
\]
This is the factorial base taken in its canonical form.  The kernel checks the
floor formula, the digit bounds $0\le d_m<m$, the recurrence
$\theta_{m+1}=(m+1)\theta_m-d_{m+1}$, the finite telescoping expansion
\[
 x=\lfloor x\rfloor+\sum_{m=2}^{N}\frac{d_m}{m!}
   +\frac{\theta_N}{N!},
\]
and the propagation rule: a zero remainder at one index forces every later
digit to vanish.  These are
\pdecl{canonicalDigit_eq_floor_mul_remainder},
\pdecl{canonicalDigit_nonneg},
\pdecl{canonicalDigit_lt_radix},
\pdecl{canonicalRemainder_recurrence},
\pdecl{factorial_expansion_partial}, and
\pdecl{zero_remainder_forces_zero_digit_tail}; they hold for every real $x$,
not only for $\Ser$.

The kernel also checks the criterion that makes the representation useful:
the digits $d_m(x)$ are eventually zero if and only if $x$ is rational
(\pdecl{exists_rat_eq_iff_eventually_zero_canonicalDigit}, in
\path{CanonicalFactorialTermination.lean}).  Specialising to $\Ser$ does not
decide whether $\Ser$ is rational.

Note what the criterion is not.  Canonical normalisation is an exact
reformulation of rationality.  It does not by itself supply an obstruction, and
a zero digit is not the same event as a zero-branch hit: the returned data
contain canonical zero digits at $m=5$ and $m=23$, while the zero-branch list
is empty through $m=100000$.

% paper-assimilation: second_layer_finite_digit_and_floor_stability_boundary sha256:75c5320d621ba7b42e0fe633783599a81ee8ad2b4b9ab08ef366dfdd555d3008
\subsection{The finite second layer and the floor-stability boundary}

There is a sharper coordinate for the actual carry sequence.  The elementary
identity
\[
 \frac1{n!-1}=\frac1{n!}+\frac1{n!(n!-1)}
\]
splits the prefix as
\[
 H_m=E_m+P_m,
 \qquad
 E_m=\sum_{2\le n\le m}\frac1{n!},
 \qquad
 P_m=\sum_{2\le n\le m}\frac1{n!(n!-1)}.
\]
For $m\ge2$, the first part is integral at factorial scale and satisfies the
exact radix recurrence $m!E_m=m((m-1)!E_{m-1})+1$.  Define the
strict-successor carry by $Z_m=mZ_{m-1}+1-b_m$.  If
\[
 f_m=\lfloor m!P_m\rfloor,
 \qquad d^{(2)}_m=f_m-mf_{m-1},
\]
then Lean proves the finite digit law
\[
 d^{(2)}_m=m-1-b_m.                                      \tag{2.1}
\]
After removing the factorial-scale integral prefix $E_m$ of $e-2$, the
strict-successor carry is exactly the defect from the maximal second-layer
digit.

Put $S_2=\sum_{n\ge2}1/(n!(n!-1))$.  The passage from $P_m$ to $S_2$ is the
delicate point.  Formula~(2.1) transfers to the canonical digit of $S_2$ only
when the floors at both scales $m!$ and $(m-1)!$ are stable.  Positivity of the
omitted tail and an upper bound below one do not suffice: a floor changes
whenever the finite prefix already lies closer to the next integer than that
tail bound.  The source therefore makes the two floor-stability hypotheses
explicit.  It also records the normalization fact that canonical factorial
digits of a rational number eventually vanish.  Hypothetical rationality of
$S$ does not imply rationality of $S_2=S-(e-2)$, since $e-2$ is irrational.
An exact-Fraction computation checks the relevant next-integer gaps through
$m=500$.  A cofinal floor-stability theorem is still needed if one wants to
drop the wrap terms.

The source already records an unconditional identity, with no floor-stability
premise.  For $m\ge3$, writing $\delta_m=m!(C-C_m)$ and
$\sigma_m=\mathbf 1_{\{\{m!C\}<\delta_m\}}$, Lean proves
\[
 b_m=m-1-d_m(C)+\sigma_m-m\sigma_{m-1}
\]
(\pdecl{factorialGapStepCarry_eq_companion_digit_law}).  Floor stability is
needed only if one discards both wrap terms.  Equality in the indicator is the
no-wrap case.

A separate exact reformulation runs through a defect automaton.  For a rational
centre recurrence $F_m=mF_{m-1}+1+\varepsilon_m-C_m$, the kernel checks that the
integer ceiling defect code equals $\lfloor m\delta_{m-1}-\varepsilon_m\rfloor$
and that $\delta_m=m\delta_{m-1}-\varepsilon_m-q_m$, with the specialisation
$\varepsilon_m=1/(m!-1)$ written out.  What is checked is the algebra of the
automaton.  Proving that the finite-sum residual centre satisfies the premise is
a separate step and is not done.

A nearby floor criterion gives a precise test for a natural approach.  Koepf
and Schmersau prove that eventual equality between the
floors of $n$ times a partial sum and $n$ times its limit forces irrationality
\cite[Theorem~1.1, p.~117]{koepf-schmersau}; their rational-term version obtains
that equality from prefix integrality at a scale $p_n$ and the strict tail
bound $a-s_n<1/(np_n)$
\cite[Theorems~2.2--2.3, pp.~119--120]{koepf-schmersau}.  For the natural
termwise clearing choice
\[
  p_n=\operatorname{lcm}\{k!-1:2\le k\le n\},
\]
the last two denominators already show the obstruction:
\[
  p_n\ge
  \frac{(n!-1)((n-1)!-1)}{n-1},
\]
because
$\gcd(n!-1,(n-1)!-1)=\gcd((n-1)!-1,n-1)\le n-1$.
For $n\ge4$, the first omitted summand $1/((n+1)!-1)$ is then already larger
than $1/(np_n)$, so this natural $p_n$ cannot satisfy their tail hypothesis.
Cancellation in the reduced prefix denominator could in principle give a
smaller scale, but proving enough cancellation is another form of the present
denominator problem.  This calculation explains exactly why their hypothesis
fails for the natural clearing denominator.

Duverney's fast-series criteria fail at a different, equally exact boundary.
His Theorem~3.1 assumes two-sided quadratic denominator growth
$cu_n^2\le u_{n+1}\le c'u_n^2$, while for $u_n=n!-1$ one has
$u_{n+1}/u_n^2\to0$ \cite[pp.~275, 285--286]{duverney}.
The all-positive specialization in Corollary~3.2 additionally requires
\[
  \sum_n\left|\frac{u_{n+1}}{u_n^2}-1\right|<\infty,
\]
whereas the summands tend to one here
\cite[Corollary~3.2, p.~287]{duverney}.  Neither criterion applies.

The sharp recent theorem of Barreto, Kang, Kim, Kova\v{c}, and Zhang has a
similarly explicit ceiling.  Its $d=1$ case proves irrationality of
$\sum_n1/a_n$ when $a_n^{1/2^n}\to\infty$, whereas
\[
  (n!-1)^{1/2^n}\longrightarrow1
\]
for the present choice $a_n=n!-1$
\cite[Theorems~2--3, pp.~2--4]{barreto-et-al}.  The proof nevertheless
identifies a useful exact condition: Mahler's elementary rationality floor is
contradicted by prefix-clearing integers $D_N$ for which the cleared positive
tails satisfy $\liminf_ND_Nr_N=0$
\cite[Lemma~8 and Proposition~12, pp.~6, 9--12]{barreto-et-al}.  The ordinary
product of the factorial-gap denominators is far too large for that estimate;
a transfer would need a low-height clearing subsequence, or enough exact
cancellation in their least common multiple.  Thus the new theorem supplies a
precise target inequality and adaptive-cutoff architecture, but not the
missing arithmetic bound.

The ordinary factorial-series direction survives more usefully.  Dividing
the strict-successor recurrence $Z_m=mZ_{m-1}+1-b_m$ by $m!$ and telescoping
gives the exact finite identity
\[
  \frac{Z_M}{M!}
  =\frac{Z_2}{2!}
   +\sum_{m=3}^{M}\frac{1-b_m}{m!}.
\]
Thus the carry defects $1-b_m$ are genuine factorial-series coefficients.
Han\v{c}l and Tijdeman give exact rationality classifications for polynomial
coefficients and finite-difference criteria for broader ordinary factorial
series \cite[Theorem~3.1 and Corollary~3.1, pp.~390--391]
{hancl-tijdeman}.  Their denominator is the cumulative linear product
$\prod_{n\le N}(an+b)$, not the individual number $N!-1$.  Applied to the
display above, the classical Cantor--Oppenheim criterion still needs
$1-b_m\ne0$ infinitely often.  Proving this is exactly the unresolved
cofinal non-unit-carry assertion.

\section{Integral channels and the LCM tax}\label{sec:channels}

Fix $d\ge2$.  The kernel proves the divisibility
\[
  (d!)^{\lfloor i/d\rfloor}\ \Big|\ i!
  \qquad(i\ge0),
\]
defines the integral channel weight $W_{d,i}$ obtained by cancelling that
factor, and checks the exact cancellation.  It also checks the consecutive
channel event
\[
  n\,W_{d,n-1}-W_{d,n}=0
  \qquad\text{whenever } d\nmid n .
\]
So the channel weight is arithmetically inert except at multiples of $d$.
The checked source statements are
\pdecl{factorial_pow_floor_dvd_factorial},
\pdecl{channelWeight_mul_denominator}, and
\pdecl{channelEvent_eq_zero_of_not_dvd}.

Let $\lambda$ be a finitely supported integer vector on indices $n\ge2$, let
$M=M(\lambda)$ be its factorial moment, and let $\chn d\lambda$ be the $d$-th
channel numerator.  The kernel checks two facts about them.

\begin{theorem}[channel congruence]\label{res:congruence}
For every finite integer support and every $d\ge2$,
\[
  \chn d\lambda\equiv M(\lambda)\pmod{d!-1}.
\]
Consequently $d!-1$ divides $M(\lambda)-\chn d\lambda$, a vanishing $d$-th
channel forces $(d!-1)\mid M(\lambda)$, and annihilating every channel
$2\le d\le D$ forces
\[
  \Lc_D=\lcm_{2\le d\le D}(d!-1)\ \Big|\ M(\lambda).
\]
\end{theorem}

\noindent On the manuscript support $n\ge2$, Lean also checks
$12\mid M(\lambda)-2\chn 2\lambda$.  Combined with $\gcd(12,\Lc_D)=1$, annihilating
channels $2$ through $D$ with $D\ge2$ forces $12\Lc_D\mid M(\lambda)$
(\pdecl{twelve_channelLCM_dvd_factorialMoment_of_channels_zero}, in
\path{ErdosProblems/Erdos68/DivisorChannelBasis.lean}).  The printed combination
of support $n\ge2$ with moment exactly $\Lc_D$ is therefore impossible.
This does not raise the parent series.

\begin{theorem}[integral normal form]\label{res:normalform}
For every finite integer support and every $d\ge2$ there is an integer $k$ with
$\chn d\lambda=M(\lambda)+(d!-1)k$.
\end{theorem}
\begin{proof}
By Theorem~\ref{res:congruence}, $\chn d\lambda-M(\lambda)$ is divisible by
$d!-1$.  The integer $k$ is that quotient.
\end{proof}

\noindent Theorem~\ref{res:normalform} is the sharper of the two for design
purposes.  It says that every zero-moment variation of the support changes the
normalised $d$-th channel contribution by an integer only.  The support-matroid
layer therefore cannot manufacture an extra fractional cancellation coordinate:
the congruence forces every normalised channel defect to be integral.

Theorem~\ref{res:congruence} imposes an arithmetic cost.  Any localiser that
kills the low channels in order to isolate the tail has moment divisible by
$\Lc_D$, and $\Lc_D$ grows faster than the tail shrinks.  The relevant
quantitative estimate is
\[
  \log \lcm_{2\le n\le N}(n!-1)\ \gg\ N^{4/3}\log N
  \qquad(N\to\infty).
\]
Here is the source-checked derivation.  Put
\[
  Q_N=\prod_{2\le n\le N}(n!-1),
  \qquad
  L_N=\lcm_{2\le n\le N}(n!-1).
\]
For an odd prime $p$, let $m_p$ count the indices $2\le n\le N$ for which
$p\mid n!-1$.  Such an index necessarily satisfies $n<p$.  The factorial
congruence multiplicity estimate of Garaev, Luca, and Shparlinski
\cite[arXiv v1, Thm.~12, p.~16]{garaev-luca-shparlinski}, applied on the
interval $1\le n\le\min(N,p-1)$, therefore gives
$m_p\ll N^{2/3}$.  (The prime $2$ divides none of these factors.)  If
\[
  E_p=\max_{2\le n\le N}v_p(n!-1),
\]
then
\[
  \log Q_N
  =\sum_p\sum_{n=2}^{N}v_p(n!-1)\log p
  \le \bigl(\max_p m_p\bigr)\sum_p E_p\log p
  \ll N^{2/3}\log L_N .
\]
On the other hand, Stirling summation gives
$\log Q_N\asymp N^2\log N$, proving the displayed lower bound for
$\log L_N$.  The cited source states the multiplicity theorem.  The lcm
corollary is derived here and has not been formalized in Lean.  It is recorded
because it quantifies the denominator growth encountered by this method.  No
later argument depends on it.

A stronger elementary bound is already available from the kernel-checked
segment inequality, without a multiplicity theorem.  Lean proves
$\prod_{n=N-k+1}^{N}(n!-1)\le L_N N^{\binom{k+1}{3}}$.  Stirling packaging of
that finite bound, with $k=\lfloor\alpha\sqrt N\rfloor$ and $\alpha=\sqrt2$,
gives the ordinary corollary
\[
  \liminf_{N\to\infty}\frac{\log L_N}{N^{3/2}\log N}\ \ge\ \frac{2\sqrt2}{3}.
\]
The constant is optimal for this coarse terminal-block estimate, not a proof
that the actual lcm exponent is $3/2$.  The exact lcm loss is the excess
repeated-layer mass $E_N=\sum\log d_n-\log L_N$, not the total pairwise gcd
mass $G_N$; those two quantities need not have the same order, and
$G_N=O(N^{3/2}\log N)$ would force almost full product growth
$\log L_N\sim\tfrac12 N^2\log N$.  The same $N^{3/2}$ lower bound holds for
$\operatorname{lcm}(n!+P(n))$ for every fixed nonzero $P\in\mathbb Z[X]$,
while $P=0$ collapses to $N!$.  The gcd identity
$\gcd(i!+P(i),j!+P(j))\mid P(j)-(j!/i!)P(i)$ is source-landed in
\path{ErdosProblems/Erdos68/PolynomialPerturbationLcm.lean}; this wave's
kernel build was deferred behind a sibling Lake owner and is not a Lean
receipt.  Lai records the printed congruence \cite[(2.5)]{lai}.  The liminf
is an ordinary proof.  None of these growth statements is a lower bound for
$\operatorname{den}(H_N)$.  In fact $L_N(S-H_N)\to\infty$, so enlarging the
full common denominator cannot make that particular positive tail small.

The primitive lcm divisibility after cofactor removal is also checked:
factorial valuations do not absorb the tax once every common cofactor divisor
has been removed.  The separate corank-one cofactor/determinant construction
for constructing such primitive kernels remains an unformalized derivation.
The checked statement here is only the lcm divisibility after cofactor removal.

\section{The prime unit translator}\label{sec:translator}

The channel tax raises an obvious question: is there any support that pays the
tax at exactly one channel?  There is, and it is explicit.

\begin{theorem}[prime unit translator]\label{res:translator}
Let $p\ge3$ be prime and take the coefficient--index pair $(p,-1)$ on the
indices $(p-1,p)$.  Then the factorial moment is $0$; every channel $d<p$
vanishes, by the exact quotient identity
$\lfloor(p-1)/d\rfloor=\lfloor p/d\rfloor$; every channel $d>p$ vanishes,
because both indices lie below $d$; and the $p$-channel numerator is exactly
$p!-1$.
\end{theorem}

\noindent Under the support convention $n\ge2$, the case $p=2$ is
inadmissible: its first index is $1$.  Omitting that index leaves moment
$-2$, not zero.  Lean also checks the identity with index $1$ admitted, but
every remote prime in the construction below may be taken at least $3$, so
the $n\ge2$ convention is retained.  For $p\ge3$ the proof supplies, at zero
cost in the moment, a unit in the $p$-channel.  Adding an integer multiple of
the translator to any
candidate kernel shifts the $p$-channel numerator by multiples of $p!-1$ and
leaves every other channel and the moment untouched.

The consequence is already uniform in the support location.  For every channel
rank and every prescribed cutoff, Lean constructs a factorial-grid kernel and
a remote prime-translator pair entirely beyond that cutoff, with all requested
low channels zero, nonzero factorial moment, and residual in $[-1/2,1/2]$;
see \pdecl{exists_remote_factorialGrid_primeTranslator_reduction}.  The rounded
residual may still equal zero, and strict nonvanishing has not been proved for
any cofinal family.  For one kernel, even strict nonvanishing only excludes
denominators dividing the associated moment.  On the compressed primitive grids
below, those moments eventually absorb every fixed denominator, so cofinal
nonintegrality is equivalent to irrationality for that family; the
nonintegrality producer itself remains open.  Section~\ref{sec:open} records
the residual target.  The ordinary proofs sit in
\path{ErdosProblems/Erdos68/CompressedPrimitiveChannelKernel.md}.

\subsection{Compressed primitive kernels}
\label{sec:compressed-kernel}
The live Cramer construction uses step $(D!)^2$.  The same channels vanish on
the smaller arithmetic progression of step $L=\operatorname{lcm}(2,\ldots,D)$.
With $i_j=r+jL$, $N=r+(D-1)L$, $a_d=(d!)^{L/d}$, $H(X)=\prod_{d=2}^D(a_d X-1)$,
and $A=\prod a_d$, the primitive integer kernel is
$\lambda_j^*=N! h_j/(A i_j!)$, with last coordinate $1$.  Its moment is
$N!\prod(1-1/a_d)$.  Every fixed positive integer eventually divides the
moments of any unbounded family of these grids.

% paper-assimilation: factorial_gap_cofinal_carry_irrationality_criterion sha256:6460cc54525e72e37fc1c20104a8845db63c53fc299397fede041a1b033cf9f7
\section{Cofinal strict-successor divisibility misses}\label{sec:plateau}

\subsection{Grid plateaux and first exit}

A second approach uses the rational grid.
Let $H$ be a partial sum and $q$ a candidate denominator.  The kernel checks
the algebraic grid threshold: writing $qH=k+r$ and $q(\Ser-H)=u$, the next
$q^{-1}$ grid point $(k+1)/q$ lies below $\Ser$ exactly when $1\le r+u$.  It
also checks the factorial plateau theorem: if $H<G\le \Ser$, if $n!G$ is
integral, and if $n!(\Ser-H)<1$, then the strict successor of $n!H$ and the
canonical floor of $n!\Ser$ are the same grid integer.

Two rigidity statements follow, both checked.  Consecutive plateau floors,
scaled by the next radix, force the canonical factorial digit to vanish.  And
any first-exit offset $\delta\in[0,2)$ with carry $b=-\lfloor\delta\rfloor$
satisfies $b\in\{0,-1\}$: the exit is rigid, with exactly two alternatives.

The checked first-crossing argument continues from the exit to a denominator
lower bound.  For a rational grid level $G$, suppose that $\tau$ is its first
crossing by the literal partial sums and write
\[
  G-H_{\tau-1}=\frac{a}{v},\qquad a,v>0.
\]
Then $v\ge\tau!-1$.  On the $-1$ exit branch this strengthens to
$v\ge\tau!(\tau!-1)$.  No coprimality hypothesis on $a$ and $v$ is required.

There is also a direct obstruction at prime indices.  Let
\[
  H_m=\sum_{2\le n\le m}\frac1{n!-1},\qquad
  Z_m=\lfloor m!H_m\rfloor+1,
\]
and let $\Delta_m$ be the distance from $(m-1)!H_{m-1}$ to its strict integer
successor.  The kernel checks, for $m\ge3$, the exact criterion
\[
  m\mid Z_m
  \quad\Longleftrightarrow\quad
  1+\frac1{m!-1}<m\Delta_m\le2+\frac1{m!-1}.
\]
If $\Ser=a/q$ with $q>0$, then for every prime $p>q$ the tail bound forces
$p\mid Z_p$.  Consequently, one exact missed prime $p$ implies $q\ge p$.
The rational implementation of $Z_p$ agrees with the real-floor definition,
and exact kernel reduction gives $11\nmid Z_{11}$.  Thus every rational
representation of $\Ser$ has denominator at least $11$.

\subsection{The lossless carry criterion}

The all-index recurrence is stronger.  Define its exact carry by
\[
  Z_m=mZ_{m-1}+1-b_m .
\]
\begin{theorem}[exact carry characterization]\label{res:carry-characterization}\label{res:strict-successor-complete-characterization}
The following conditions are equivalent:
\[
  \Ser\notin\mathbb Q,
  \qquad
  (\forall B)(\exists m>B)\ b_m\ne1,
  \qquad
  (\forall B)(\exists m>B)\ m\nmid Z_m.
\]
Equivalently, $\Ser$ is rational precisely when $b_m=1$ for all sufficiently
large $m$.
\end{theorem}

\noindent
If $\Ser=a/q$ and $m-1\ge q$, the plateau theorem identifies
\[
  Z_{m-1}=\frac{(m-1)!}{q}\,a,\qquad
  Z_m=\frac{m!}{q}\,a=mZ_{m-1},
\]
so necessarily $b_m=1$.  Hence one exact non-unit carry at index $m$ forces
$q\nmid(m-1)!$, and therefore also $q\ge m$.  Conversely, if $b_m=1$
eventually, then $Z_m/m!$ is eventually constant.  The checked one-cell bound
\[
  H_m<\frac{Z_m}{m!}\le H_m+\frac1{m!}
\]
and the exact tail estimate show that $Z_m/m!\to\Ser$; hence that eventual
constant is $\Ser$ and is rational.  Theorem~\ref{res:carry-characterization}
therefore reduces the original irrationality question to proving the
existence of arbitrarily late non-unit carries.  Exact
rational normalization gives
\[
  60\nmid Z_{60},\qquad 64\nmid Z_{64},\qquad 67\nmid Z_{67}.
\]
At $m=60$ the recurrence also proves $b_{60}\ne1$, and hence $q\ge60$.
Since $67$ is prime, the prime-miss theorem applied at $67$ gives the stronger
checked bound
\[
  \Ser=\frac aq,\ q>0 \quad\Longrightarrow\quad q\ge67.
\]

There is also a sufficient arithmetic specialization at doubled prime indices.  For
every odd prime $p$, the kernel now specializes the strict-successor
prime-power criterion to the literal prefixes:
\[
  p^2\mid Z_{2p}
  \quad\Longleftrightarrow\quad
  \bigl(b_{2p}=1\ \text{and}\ p\mid Z_{2p-1}\bigr)
  \ \text{or}\
  \bigl(b_{2p}=1+p\ \text{and}\ p\mid 2Z_{2p-1}-1\bigr).
\]
Consequently, failure of both displayed branches for a cofinal family of odd
primes proves $\Ser$ irrational.  This is a sharper two-stage target than a
bare square nondivisibility assertion: it exposes separately the only two
carry values and predecessor residues that can occur.  The cofinal failure of
both branches remains unproved.

A separate exact GMP integer computation certifies all $299998$
carry cells for $3\le m\le300000$.  Its unit carries occur exactly at
\[
  52,\ 591,\ 1030,\ 1407,\ 1438,\ 2164,\ 4258,\ 10991,\ 21236,
\]
so no further unit carry occurs through $m=300000$, where
$b_{300000}\ne1$.  Applying the checked divisibility theorem gives
\[
  \Ser=\frac aq,\ q>0 \quad\Longrightarrow\quad q\nmid299999!.
\]
In particular $q\ge300000$, but the divisibility statement is stronger: it
excludes every divisor of $299999!$, including divisors larger than $300000$.
It does not exclude all smooth denominators: a sufficiently high prime power
can fail to divide $299999!$ even when its prime is small.
The computation uses no floating-point arithmetic; its canonical data,
Python driver, and GMP backend are hash-bound in the companion research
packet.  This remains a finite exclusion.  Irrationality requires the same
conclusion beyond every cutoff.

An independent continued-fraction calculation is recorded, with its source
and receipt, in the research packet.  At an 80000-bit exact bracket it forces
23449 common partial quotients and verifies
the penultimate convergent's separation from $\Ser$, yielding
$q\ge2^{39990}>10^{12038}$ by best approximation of the second kind.  This
is also a finite exclusion.  It bounds the size of $q$, whereas the carry
calculation restricts its prime-power multiplicities.
The same enclosure admits an ordinary interval-optimal extraction from the
first differing complete quotients, of the form $Q_*=k q_j+q_{j-1}$.  On an
open remaining interval the integer is $k=\lfloor\alpha\rfloor+1$; on a closed
interval it is $k=\lceil\alpha\rceil$.  After the prefix $[0;1,2]$, the closed
complete-quotient interval $[49,51]$ maps to $[103/154,99/148]$, whose minimal
denominator is the included endpoint rather than the open-interval candidate
$k=50$.  This is a convention correction, not a new enclosure and not a new
numerical floor.  The live lead remains the last common convergent
$q\ge2^{39990}>10^{12038}$; the remaining complete-quotient interval is not
stored in the receipt, and the $80000$-bit scan is not rerun.

\subsection{The adjacent-unit window reduces to a universal bound}
\label{sec:adjacent-unit-no-go}

Consider an attempt to exclude two consecutive unit carries by extracting a
prime-power or quotient-gcd obstruction from their cleared width-one window.
The exact two-step recurrence reduces this proposed obstruction to a universal
bound.  Let $U_k$ and $V_k$
be the positive numerator and denominator of the reduced predecessor gap, and
let $G_k>0$ be the exact transition normalizer.  For $m\ge3$, the two unit
carries are equivalent to an integer offset $\Omega_m$ satisfying
\[
  0<\Omega_m\le D_m.
\]
The window denominator telescopes independently of the carry values:
\[
  D_m=V_{m+2}G_{m+1}G_m
     =V_m(m!-1)((m+1)!-1).
\]
More generally, before imposing either carry value, the cleared offset has the
exact mixed-radix factorisation
\[
 \Omega_m=U_{m+2}G_{m+1}G_m
   +D_m\bigl((m+1)b_m+b_{m+1}-(m+2)\bigr).
\]
The second summand records the complete two-carry defect.  This identity
identifies the Archimedean window and
the reduced-denominator dynamics as two descriptions of the same integer.
Under the adjacent-unit assumption, the offset has the matching
factorisation
\[
  \Omega_m=U_{m+2}G_{m+1}G_m.
\]
Cancelling the common positive normalizer therefore reduces the complete
window to
\[
  0<U_{m+2}\le V_{m+2},
\]
which is the universal bound for a reduced strict-successor gap.

Consequently, after the two unit carries are imposed, this particular cleared
window is the universal reduced-fraction bound rewritten through a common
factor.  Further arithmetic restrictions on the actual predecessor state would
be additional information; the telescope does not show that every such
restriction is redundant.

% paper-assimilation: nearest_integer_linear_form_divisibility_boundary
\subsection{Why nearest-integer nonvanishing is only a denominator screen}

The nearest-integer residual used by a rank-two channel construction has a
precise rational meaning, but it is not by itself a small-linear-form
contradiction.  Write a hypothetical value as $\Ser=a/q$ in lowest terms and,
for integers $N,K$, put
\[
  \delta= N\Ser+K-\operatorname{nint}(N\Ser+K),
  \qquad -\tfrac12\leq\delta\leq\tfrac12 .
\]
Then $q\delta$ is integral.  Consequently
\[
  \delta=0\quad\Longleftrightarrow\quad q\mid N,
  \qquad
  \delta\ne0\quad\Longrightarrow\quad |\delta|\geq\frac1q.
\]
The universal nearest-integer bound therefore remains perfectly compatible
with rationality: for $\Ser=1/q$ and $N=1$, the residual is nonzero while it
still lies in $[-1/2,1/2]$.  Its strict nonvanishing can only screen the
denominator through $q\nmid N$; it supplies no irrationality conclusion
without a stronger, denominator-sensitive comparison.

For the factorial-gap rank-two construction, the checked Cramer identity
identifies $N$ with the appended determinant $\det A$, while the prime
translator changes only an integral coordinate.  Thus a nonzero rounded
residual says exactly that a hypothetical denominator does not divide
$\det A$.  To prove irrationality by this method for one kernel, one would need such
nonvanishing cofinally.  On the compressed primitive grids of
\S\ref{sec:compressed-kernel} the primitive moments eventually absorb every
fixed denominator, so an extra bound below $1/q$ is unnecessary for that
family; cofinal nonintegrality remains open.  This does not revive the
refuted certificate $0<|N|<\gcd(\mathrm{minors})$.

There is also a checked finite peeling identity.  For $x\ne0,1$ and $K\ge0$,
\[
  \frac1{x-1}
  =\sum_{j=1}^{K}\frac1{x^j}
   +\frac1{x^K(x-1)}.
\]
When $x=k!$ and a chosen factorial scale is divisible by $(k!)^K$, the scaled
finite sum is integral and only the last term retains the factor $k!-1$ in
its denominator.  The identity isolates one residual fraction before exact
bounding; it does not yet give a cofinal family of nonzero residuals.

One proposed strengthening is false and is recorded as such: the divisibility
$(m!-1)\mid\operatorname{den}(V_m)$ fails at the reported strict events $m=52$
and $m=591$.  Only the Archimedean first-crossing lower bound survives.

The carry equivalence states the full remaining problem.  Wilson factors,
doubled-prime branches, congruence tests, nearest-integer residuals, and the
finite GMP computation give ways to study individual indices or exclude
denominators.  Irrationality still requires the quantifier
$(\forall B)(\exists m>B)$.

% paper-assimilation: factorial_gap_prime_pole_principal_residue sha256:6a0d40220022d5a36d5f40e488f7850374c6f7994d12f37ffa0d8d407fc52ad6
\section{Finite prime-pole cancellation}
\label{sec:prime-pole}

The lcm of a finite prefix records the largest prime-power pole, but it does
not say whether that pole survives when the summands are added.  This
cancellation is now exact.  Put
\[
 L_M=\mathop{\rm lcm}_{2\le n\le M}(n!-1),
 \qquad
 A_M=\sum_{2\le n\le M}\frac{L_M}{n!-1},
\]
so that $H_M=A_M/L_M$ before reduction.  Fix a prime $q$, and suppose $e\ge1$
is the largest exponent for which $q^e$ divides at least one displayed
denominator.  Let
\[
 I_{q,e}(M)=\{n\in[2,M]:v_q(n!-1)=e\}
\]
and define its principal residue
\[
 R_{q,e}(M)=
 \sum_{n\in I_{q,e}(M)}
 \left(\frac{n!-1}{q^e}\right)^{-1}\pmod q .
\]
All terms of lower $q$-adic order vanish after multiplication by $L_M$ and
reduction modulo $q$.  Lean proves the surviving identity
\[
 A_M\equiv \frac{L_M}{q^e}\,R_{q,e}(M)\pmod q.           \tag{6.1}
\]
Maximality of $e$ makes $L_M/q^e$ a unit modulo $q$.  Consequently
\[
 v_q\!\left(\operatorname{den}H_M\right)=e
 \quad\Longleftrightarrow\quad
 R_{q,e}(M)\ne0\pmod q.                                  \tag{6.2}
\]
This is the exact finite prime-pole criterion: support and multiplicity are
compressed into one residue-weighted sum, with no loss between the displayed
factorial gaps and the reduced prefix.  For a general finite sum of
reciprocals the same first-layer criterion is a direct consequence of
Louwsma and Martino's valuation formula \cite[Lemma~4.1]{louwsma-martino},
which retains the full surviving exponent rather than only the top layer.
The short argument above is included to fix the factorial-gap normalisation;
the novelty, such as it is, lies in the factorial examples and the
downstream consumers, not in the general residue identity.

The distinction is arithmetically real.  At $M=q-1$, Wilson's theorem always
provides the endpoint hit $q\mid(q-2)!-1$, but it does not isolate that hit or
prevent cancellation at the maximal layer.  The first exact cancellation
found by the local scan is
\[
 q=139,\qquad I_{139,1}(138)=\{69,122,137\},
\]
with maximal-hit cofactors $6,49,73$ and
\[
 6^{-1}+49^{-1}+73^{-1}=0\pmod{139}.
\]
The second is
\[
 q=2593,\qquad I_{2593,1}(2592)=\{349,2243,2591\},
\]
with cofactors $1508,1566,1678$ and reciprocal sum zero modulo $2593$.
Lean checks both displayed reciprocal equalities by exact computation.  A
separate local modular scan through $q\le3000$ reproduces precisely these two
examples and all the displayed data.  The larger external scan through
$200000$ has not been locally regenerated and is not used as theorem
authority here.

Equations~(6.1)--(6.2) have two immediate consequences.  First, a large
prime factor or Wilson endpoint hit is not automatically a pole of the
reduced prefix; one must control the entire maximal-hit residue.  Singleton
maximal support is the easy noncancelling case, since its cofactor is a unit.
Second,
incidence bounds that count hits but ignore their inverse cofactor weights can
at best control the size of $I_{q,e}(M)$, not the cancellation that decides
survival.  A proof along these lines would need a cofinal family of primes and
prefixes for which this weighted residue is nonzero, followed by an
independent Archimedean argument linking the surviving pole to the
strict-successor interval.  Equations~(6.1)--(6.2) settle the local arithmetic
part of that problem.

% paper-assimilation: four_step_affine_defect_rigidity sha256:e3ed6a89baf7aaaf81d258bf0ddeddcbcebc721701bb6fa47d124718aa0659d7
\section{A four-step defect escape criterion}
\label{sec:four-step-defect}

The following recurrence gives another necessary condition.  Let $e_n$ and
$\eta_n$ be integer sequences related by
\[
  \eta_n=(n+1)e_n-e_{n+1},
  \qquad Q(n)=n^2+2.
\]
Lean checks the following four-step alternative.  If $e_n>0$ and
$\eta_{n+j}<n+j$ for $j=0,1,2,3$, then
\[
  e_{n+4}>Q(n+4).
\]
Applying the same calculation to $-e$ and $-\eta$ gives the signed form: if
$e_n\ne0$ and $|\eta_{n+j}|<n+j$ through the same four positions, then
$|e_{n+4}|>Q(n+4)$.  Consequently an eventual quadratic window
$|e_n|\le Q(n)$ together with eventual sublinear defects forces $e_n=0$
eventually.  In the one-sided factorial-neighbour specialization, cofinally
many nonzero errors inside that window force cofinally many indices with
$\eta_k\ge k$.

The proof is elementary but the integer gap matters.  From $e_n\ge1$, rewrite
$e_{r+1}=(r+1)e_r-\eta_r$ four times and insert the largest permitted defect
at each step.  The resulting integral lower bound lies beyond the quadratic
window.  Sign normalization supplies the two-sided statement.

The checked recurrence quantifies the growth forced on a factorial-neighbour
approximation.  Applying it to the present series would require a reciprocal
weighted transform connecting these errors to $\Ser$, a proof of the stated
window for that transform, and cofinal defects for the actual strict
successors.  None of these three inputs is presently available.  The open
step remains to prove $m\nmid Z_m$ beyond every cutoff.

\section{Weighted projection rigidity}\label{sec:projection}

The third checked layer converts modular disagreement into exclusion.

\begin{theorem}[projection rigidity]\label{res:projection}
Suppose the natural endpoint numerator $Z$ is congruent to a weighted numerator
$T$ modulo $R$.  Then every divisor $Q$ of $R$ with $Z\le B<Q$ satisfies
$T\bmod Q=Z$.  Consequently two divisors $Q_1,Q_2>B$ with unequal projected
residues exclude any such bounded endpoint; and unequal projections force
$\min(Q_1,Q_2)\le T$.
\end{theorem}

\noindent The leave-one-out specialisation $Q=R/r$ is also checked.  More
generally, let $a,b$ divide $R$ and take the complementary projection moduli
$R/a$ and $R/b$.  Lean checks the same quotient cancellation and
collision-cap comparison for these factor projections.  If $\gcd(a,b)=1$,
then $\lcm(R/a,R/b)=R$, and the resulting branch-free factor-pair floor is at
most the global complementary residue.  The factors $a,b$ may both divide one
private quotient.  Thus the checked reduction needs no analytic input and no
pair of distinct denominator indices, only suitable factors whose projections
or factor-pair floor satisfy the stated bound.  This is checked in
\pdecl{factorialBlock_factorProjection_lcm_eq_privateModulus} and
\pdecl{factorialBlock_complementaryFactorPairFloor_le_global}.

Lean also verifies the construction for the literal factorial block.  It
builds the collision core $C$, private quotients $r_i$, private
modulus $R$, and weighted numerator $T$ for the actual denominators $i!-1$,
and proves both the endpoint congruence modulo $R$ and the required
coprimality.  Moreover, if $m!-1$ has canonical large prefix-private primes,
then their complete prime-power product divides the single quotient owned by
$m$ on the tailored block with parameter $\lfloor m/2\rfloor+1$, hence divides
that block's $R$.

The collision core itself now has an exact incremental law.  For the positive
factorial-gap denominators, adjoining $d_a$ to an old finite family $S$ gives
\[
 C(S\cup\{a\})
 =\lcm\!\left(C(S),\,
     \gcd\!\left(d_a,\lcm_{j\in S}d_j\right)\right).
\]
Indeed, finite-family gcd--lcm distributivity collapses the lcm of all
pairwise gcds against $d_a$ to this single gcd.  The same formula holds after
adjoining the distinguished base; see
\pdecl{pairwiseCollisionCore_insert_gcd_lcm} and
\pdecl{collisionCore_insert_gcd_lcm}.  Thus each step needs only the old
denominator lcm and the old collision core, with no pairwise rescan.

There is now also an exact product--lcm bound.  If
$\widetilde C(S)$ denotes the collision core after cancelling a positive
distinguished base, while
$L(S)=\lcm_{j\in S}d_j$ and $P(S)=\prod_{j\in S}d_j$, then Lean proves
\[
 \widetilde C(S)L(S)\mid P(S),
 \qquad\text{hence}\qquad
 \widetilde C(S)\le \frac{P(S)}{L(S)}.
\]
See
\pdecl{collisionCore_div_base_mul_denominatorLcm_dvd_denominatorProd}.
For the actual factorial block this specializes to
\[
 \operatorname{factorialBlockNormalizedCollisionCore}(p)
 \le
 \frac{\displaystyle\prod_{n\in I_p}(n!-1)}
      {\displaystyle\lcm_{n\in I_p}(n!-1)},
\]
where $I_p$ is the block index set; see
\pdecl{factorialBlockNormalizedCollisionCore_le_gapProd_div_gapLcm}.
This inequality converts lower estimates for the factorial-gap lcm into upper
estimates for the normalized collision core.  Its use in an irrationality
proof would still require cofinal estimates strong enough at the selected
private factor and factorial scale.  A
fixed-modulus hit count alone does not supply such control.

The distinguished-base cancellation is now exact prime by prime.  Writing
$B=(p-1)!$ and $C$ for the unnormalised factorial-block collision core,
\[
 \widetilde C_p=\frac{\lcm(B,C)}{B}
 =\frac{C}{\gcd(B,C)},\qquad
 v_q(\widetilde C_p)=v_q(C)-\min\{v_q(B),v_q(C)\}.
\]
See \pdecl{collisionCore_div_base_eq_pairwiseCollisionCore_div_gcd},
\pdecl{collisionCore_div_base_factorization}, and
\pdecl{factorialBlockNormalizedCollisionCore_factorization}.  Consequently
$q^e\mid\widetilde C_p$ exactly when the pairwise core carries
$q^{e+v_q(B)}$; in the factorial block this forces two distinct gaps to be
divisible by that higher power.  Lean moreover proves the sharp surviving
valuation cap
\[
 v_q(\widetilde C_p)+v_q((p-1)!)<q.
\]
Thus every support prime satisfies
$p-1<q(q-1)<q^2$, and, whenever $k(k-1)\le p-1$,
$\widetilde C_p$ is coprime to $k!$; see
\pdecl{factorialBlock_normalizedCollision_factorization_add_base_lt_prime},
\pdecl{factorialBlock_prime_sq_gt_pred_of_dvd_normalizedCollisionCore}, and
\pdecl{factorialBlockNormalizedCollisionCore_coprime_factorial_of_mul_pred_le}.
This removes every factorial channel below the moving square-root cutoff.
The aggregate product of the remaining large prime powers at the selected
quotient, as well as the complementary projections and residues, still needs
cofinal control.

For collision estimates that already provide an upper-half hit, no exponent
is lost to normalization.  If $q$ divides a displayed factorial gap at some
$n\ge p$, then $q\nmid(p-1)!$, and Lean proves for every $e>0$ that
\[
 q^e\mid\widetilde C_p\quad\Longleftrightarrow\quad q^e\mid C_p.
\]
See \pdecl{factorialBlockPrime_not_dvd_base_of_upper_hit} and
\pdecl{factorialBlock_primePower_dvd_normalizedCollisionCore_iff_of_upper_hit}.
Combined with the two-hit theorem, this identifies every complete normalized
upper-hit contribution with repeated full-power load in two distinct
displayed gaps.  The remaining arithmetic task is to aggregate those moving
loads strongly enough for the normalized collision cap; this equivalence
does not provide that estimate or the complementary-residue bound.

This bridge has an exact incidence-count form.  For an upper-hit prime $q$
and every $e>0$, Lean proves
\[
 q^e\mid\widetilde C_p
 \quad\Longleftrightarrow\quad
 1<\#\{i\in I_p:q^e\mid i!-1\}.
\]
See
\pdecl{factorialBlock_primePower_dvd_normalizedCollisionCore_iff_one_lt_hitCount}.
Hence a source estimate giving at most one $q^e$-hit deletes that exponent
from the normalized core and yields
$v_q(\widetilde C_p)<e$; see
\pdecl{factorialBlock_normalizedCollisionCore_factorization_lt_of_hitCount_le_one}.
The remaining problem is genuinely aggregate: obtain sufficiently uniform
incidence bounds over all moving support primes and exponents, multiply the
surviving valuation contributions, and still close the complementary-residue
coordinate.

The local aggregation is now exact.  For every upper-hit prime $q$, Lean
proves
\[
 v_q(\widetilde C_p)
 =
 \#\left\{e\in[1,q-1]:
   1<\#\{i\in I_p:q^e\mid i!-1\}\right\};
\]
see
\pdecl{factorialBlock_normalizedCollisionCore_factorization_eq_repeatedHitLayerCount}.
There is therefore no additional valuation loss between prime-power
incidence estimates and the complete local collision exponent.  The open
step is to bound these layer counts uniformly as $q$ and $p$ move, then
control the product over all surviving primes strongly enough for the
normalized collision cap; this theorem does not supply that global estimate.

The same local load now has a distance-sensitive witness.  Put
$B=(p-1)!$.  If $q$ is prime, $e>0$, and
$q^e\mid\widetilde C_p$, Lean produces $i<j$ in $I_p$ such that
\[
 q^{e+v_q(B)}\mid i!-1,\qquad
 q^{e+v_q(B)}\mid j!-1,\qquad
 q^{e+v_q(B)}\le j^{\,j-i}.
\]
See
\pdecl{exists_factorialBlock_hitPair_with_normalizedPrimePower_le_gapPow}.
Consequently, if
\[
 (2p-1)^d<q^{e+v_q(B)},
\]
then some such two hits satisfy $d<j-i$; see
\pdecl{exists_factorialBlock_hitPair_distance_gt_of_endpointPow_lt_normalizedPrimePower}.
The spacing hypothesis in that reduction is now discharged internally.
If $q$ is prime, then any two $q^e$-hits $i<j$ satisfy
$e<j-i$, without an endpoint or large-prime hypothesis; see
\pdecl{factorialBlock_prime_hitPair_distance_gt_exponent}.  The point is
that $q\mid j!-1$ already forces $j<q$, while the preceding gap-power
inequality converts this automatic size relation into strict separation.
Consequently Lean proves the global primewise diameter ceiling
\[
 q\mid\widetilde C_p
 \quad\Longrightarrow\quad
 v_q(\widetilde C_p)+v_q((p-1)!)<2p-3
\]
for every prime $q$ and $p\ge2$; see
\pdecl{factorialBlock_normalizedCollision_factorization_add_base_lt_blockDiameter}.
The exponent-level version is
\pdecl{factorialBlock_normalizedCollision_exponent_add_base_factorization_lt_blockDiameter}.
Thus the earlier endpoint-prime estimate is a special case, and even primes
already present in the normalization base pay for their base valuation
inside the same block-diameter budget.  This still does not control how many
collision primes occur or the product of their bounded powers; those global
estimates, together with the complementary-residue bound, remain open.

The pairwise statement is stronger than the selected-witness form used in
that proof.  For arbitrary $q,e$ and any displayed hits $i<j$, Lean proves
\[
 q^e\mid(i!-1),\quad q^e\mid(j!-1)
 \quad\Longrightarrow\quad q^e\le j^{\,j-i};
\]
see \pdecl{factorialBlock_primePower_le_gapPow_of_two_hits}.  Consequently,
when $q$ is prime and $e>0$, every two $q^e$-hits in the block satisfy
$e<j-i$, so the condition holds for every pair and one chosen pair does not
suffice; see
\pdecl{factorialBlock_prime_hitPair_distance_gt_exponent}.
Every prime-power hit layer is therefore an $e$-separated subset of the block.
Lean now proves the finite cardinality corollary itself:
\[
 (e+1)\#\{i\in I_p:q^e\mid i!-1\}\le 2p+e-2
\]
for $p\ge2$, prime $q$, and $e>0$; see
\pdecl{factorialBlock_primePowerHitCount_mul_succ_le}.
The remaining task is to combine this packing estimate with the exact
repeated-layer valuation identity, sum the
prime-power weights over all moving collision primes, and prove a global
product bound strong enough for the normalized collision cap.

For an endpoint prime carrying one upper-half hit, Lean now performs the
first combination exactly.  If $p\ge2$, $q$ is prime, and $2p-1<q$, then
\[
 v_q(\widetilde C_p)
 =
 \#\left\{e\in[1,2p-4]:
   1<\#\{i\in I_p:q^e\mid i!-1\}\right\};
\]
see
\pdecl{factorialBlock_normalizedCollisionCore_factorization_eq_truncatedRepeatedHitLayerCount}.
Thus every repeated-hit layer outside the block-diameter window has been
removed from the exact local valuation formula.  The remaining estimate is
still global and weighted: these truncated layer counts must be aggregated
over the moving endpoint primes strongly enough to bound their complete
prime-power product, and the independent complementary-residue coordinate
remains open.

The endpoint incidence criterion itself no longer needs a selected
upper-half anchor.  For every prime $q>2p-1$ and $e>0$, Lean proves
\[
 q^e\mid\widetilde C_p
 \quad\Longleftrightarrow\quad
 1<\#\{i\in I_p:q^e\mid i!-1\};
\]
see
\pdecl{factorialBlock_primePower_dvd_normalizedCollisionCore_iff_one_lt_hitCount_of_endpoint_lt_base}.
Thus an at-most-one $q^e$ incidence estimate forces
$v_q(\widetilde C_p)<e$ without first choosing an upper hit; see
\pdecl{factorialBlock_normalizedCollisionCore_factorization_lt_of_endpoint_lt_base_of_hitCount_le_one}.
At $e=2$ this gives the conditional squarefree conclusion
$v_q(\widetilde C_p)\le1$; see
\pdecl{factorialBlock_normalizedCollisionCore_factorization_le_one_of_endpoint_lt_base_of_primeSqHitCount_le_one}.
More generally the endpoint inequality can be replaced by the exact condition
$q\nmid(p-1)!$.  For every such prime and every $e>0$, Lean proves the same
hit-count equivalence; see
\pdecl{factorialBlock_primePower_dvd_normalizedCollisionCore_iff_one_lt_hitCount_of_not_dvd_base}.
The at-most-one hypothesis cuts the normalized valuation below $e$, and for
$e=2$ it gives conditional squarefreeness; see
\pdecl{factorialBlock_normalizedCollisionCore_factorization_lt_of_not_dvd_base_of_hitCount_le_one}
and
\pdecl{factorialBlock_normalizedCollisionCore_factorization_le_one_of_not_dvd_base_of_primeSqHitCount_le_one}.
Every prime $q\ge p$ is absent from $(p-1)!$, so this covers the entire moving
prime range at and above the block parameter.  The missing inputs are a
uniform prime-square incidence theorem and a global weighted product estimate.
The squarefreeness premise is not proved.  An exhaustive modular scan through
$q\le2{,}000{,}000$ and $n\le240$ found four individual square hits and no
prime with two such hits.  Separately, all $498{,}501$ pairs
$2\le a<b\le1000$ have squarefree $\gcd(a!-1,b!-1)$, and the aggregate
squarefree-collision scan through $p=499$ stays below $0.374$ of the
upper-descending-factorial logarithmic scale.  These are finite exact
computations.  They do not provide an asymptotic incidence bound.

Cofinal prefix-private support itself is unconditional.  Given any cutoff
$B$, Lean chooses a prime $q\ge B!+5$, uses Wilson's theorem to obtain
$q\mid(q-2)!-1$, and takes the least factorial-gap hit $m$ of $q$.  If
$m\le B$, then $q\le m!-1\le B!$, a contradiction.  Hence $m>B$; see
\pdecl{cofinal_prefixPrivate_factorialGap_hits}.  A finite variant compares the
product of a chosen set of primes, each at least $5$, with
\[
  \prod_{2\le k\le B}(k!-1).
\]
If the prime product is larger, at least one chosen prime has no hit through
$B$, while Wilson still bounds its least hit by $q-2$; see
\pdecl{exists_late_prefixPrivate_factorialGap_hit_of_primeProduct_lt}.
These statements supply private factors, but they do not prove either scale
estimate below.  In particular, the unconditional construction gives no
useful upper bound for $q$ in terms of its least hit $m$.

Wilson reflection also limits what can be inferred from a prime factor merely
because it is linear in a later index.  If $n$ is odd, $n<q$, and
$q\mid n!-1$, then $q\mid(q-n-1)!-1$.  When both indices lie in the same block
and the reflected hit is earlier, equivalently $q<2n+1$, this repeated hit
survives predecessor-factorial normalization and its full-block incidence
count exceeds one; see
\pdecl{prime_dvd_reflected_factorialGap_of_odd} and
\pdecl{prime_dvd_factorialBlockNormalizedCollisionCore_of_odd_reflection} and
\pdecl{one_lt_factorialBlockPrimeHitCount_of_odd_reflection}.  Thus a
linear-size divisor need not be private.  The divisibility identity itself
does not automatically supply a second admissible series index: it gives a
second hit in the factorial-gap family only when $n'=q-n-1$ satisfies
$n'\ge2$ and $n'\ne n$, and a collision in a specified block only when both
indices belong to that block.  The pair $(q,n)=(23,11)$ is a fixed point of
the reflection, while $(q,n)=(5,3)$ reflects to $n'=1$.  Lean's collision
consumers already carry the membership and ordering hypotheses; the
restriction is on the prose inference, not on the identity.

The point follows directly from Stewart's stated theorem.  For every
$\varepsilon>0$ there are infinitely
many odd $n$ whose least prime factor $q$ of $n!-1$ satisfies
\[
 n<q<
 \left(\frac{\sqrt{145}-1}{8}+\varepsilon\right)n;
\]
the printed text explicitly transfers estimate~(9) from $n!+1$ to $n!-1$
\cite[p.~464]{stewart2004}.  Wilson reflection then supplies the earlier hit
$q\mid(q-n-1)!-1$.  The source controls $q$ relative to the later index $n$,
but it does not control $q$ relative to the private first-hit index $m$.
Accordingly it contributes to the collision-core estimate; it supplies
neither a private-anchor estimate nor a global product estimate.

The strongest checked factor-level consumers can be stated without the
earlier two-owner architecture.

\begin{theorem}[moving private-factor scale split]
\label{res:moving-factor-scale-split}
For every cutoff $B$, suppose there are $m\geq4$ and a canonical large
prefix-private prime $q\mid m!-1$ such that, for
$p=\lfloor m/2\rfloor+1>B$,
\[
 (2p+1)L_p<2p^2(2p-1)!\rho_p,
 \qquad
 (2p+1)C_pq<2p^2(2p-1)!.
\]
Then $\Ser$ is irrational.
\end{theorem}

\begin{theorem}[split-factor normalized-collision criterion]
\label{res:split-factor-normalized-collision}
Suppose that cofinally many prime block parameters $p$ admit factors
$a,b\mid R_p$, with $R_p>1$, whose complementary projections modulo
$R_p/a$ and $R_p/b$ disagree and for which
\[
 (2p+1)\widetilde C_p\max\{a,b\}
   <\frac{2p^2(2p-1)!}{(p-1)!}.
\]
Then $\Ser$ is irrational.  The factors $a$ and $b$ may lie in one moving
private owner quotient; they need not come from distinct denominator indices.
\end{theorem}

In fact one selected prime $q$ already furnishes the exact
coprime factor pair $(1,q)$: its projection moduli are $R$ and $R/q$, whose
least common multiple is $R$.  Thus no second selected prime is needed; see
\pdecl{factorialGapLargePrefixPrivatePowerModulus_dvd_tailoredBlockPrivateQuotient}
and
\pdecl{factorialGapLargePrefixPrivatePrimes_unit_factor_pair_tailoredBlock}.
There is no hidden equality/disagreement branch in this specialization.
Writing $\rho$ for the global complementary residue, Lean proves that the
unit-pair floor is exactly
\[
  \min\{\rho,R/q\};
\]
see \pdecl{factorialBlock_unitFactorPairFloor_eq_min}.  Consequently the
remaining factor-pair scale comparison must simultaneously beat the global
complementary-residue coordinate and the local $R/q$ coordinate.  After using
$L=CR$, the latter is precisely the collision-cap comparison with the selected
factor $q$, while the former is the global complementary-residue lower bound.
Lean records this as the exact equivalence
\[
 (2p+1)L < 2p^2(2p-1)!\min\{\rho,R/q\}
 \quad\Longleftrightarrow\quad
 \begin{cases}
  (2p+1)L < 2p^2(2p-1)!\rho,\\
  (2p+1)Cq < 2p^2(2p-1)!,
 \end{cases}
\]
see \pdecl{factorialBlock_unitFactorPairFloor_scale_iff}.  Thus the
unit-factor reduction has no opaque floor premise left: the two surviving
arithmetic estimates are exposed independently and neither follows merely
from coprimality.  The global half is exactly the joint logarithmic budget
\[
 \log\widetilde C_p+\log(1/\delta_p)
 <\log\frac{2p^2(2p-1)!}{(2p+1)(p-1)!}
 =p\log p+(2\log2-1)p+O(\log p),
\]
with $\delta_p=\rho_p/R_p$.  Both losses must be controlled on the same
cofinal family; a bound for either term alone is insufficient, and
$\gcd(T_p,R_p)=1$ supplies only $\delta_p\ge1/R_p$.
The resulting irrationality criterion works for every natural block parameter
at least three, including composite parameters.  Wilson's theorem supplies
cofinal prefix-private factors without analytic input.  The stronger
source-backed large-prime selection supplies a positive-density family and a
linear lower bound for $q$ relative to the original hit.  To apply the
criterion, one must still prove both the global complementary-residue bound
and the local collision-core bound cofinally.  These two requirements are
packaged by
\pdecl{irrational_factorialGapSeries_of_cofinal_largePrefixPrivate_unitScaleSplit}.

\section{Obstructions to several proposed arguments}\label{sec:nogo}

The following calculations show why several natural arguments cannot prove
irrationality in their present form.

\begin{itemize}
\item Residue vectors, their recurrences, and window widths admit synthetic
  all-hit blocks.  They cannot prove irrationality on their own.
\item Known pointwise prime congruences, prime-dilation congruences, parity, and
  the exact prime coefficient formula admit a synthetic rational countermodel.
  A congruence family that a rational number could also satisfy decides nothing.
\item Wilson quotients, harmonic sums, $p$-adic gamma identities, and factorial
  residues do not control the required Archimedean floor without an additional
  coupling theorem.  Every prime-window test factors into a sharp Archimedean
  strict-ceiling condition and a modular divisibility condition, and the missing
  ingredient is a theorem coupling the Archimedean and modular
  conditions.
\item Wilson-endpoint support or hit-count bounds alone do not prove that a
  prime power survives reduction of a factorial-gap prefix.  The
  principal-residue formula of \S\ref{sec:prime-pole} is exact, and the scan's
  cancellations at $q=139$ and $q=2593$ show that the reciprocal cofactor
  weights are indispensable.  Singleton maximal support, by contrast, cannot
  cancel.
\item Fixed-denominator scalar canonical-product localisers and rank-saturated
  consecutive-jet Hermite--Pad\'e systems pay the full factorial-gap denominator.
  For $E(z)=\prod_{n\ge2}(1-z/n!)$ the genus-zero product satisfies
  $-E'(1)/E(1)=\Ser$, and the natural scalar linear form carries the coefficient
  $Q_N=\prod_{2\le n\le N}(n!-1)$, for which $Q_N$ times the tail diverges.
  Exact first-order interpolation, scalar residue weighting, the natural
  Wronskian, and rank-saturated consecutive jets all reassemble the same
  prohibitive denominator.
\item The support-matroid layer cannot create an additional fractional
  cancellation coordinate (\S\ref{sec:channels}), and factorial valuations
  cannot absorb the channel LCM tax.
\item A fixed pair of low-index private owners cannot yield a cofinal
  projection argument.  If the owner index $n$ is fixed and $p>n!-1$, then
  $n!-1\mid(p-1)!$, so its private quotient in the factorial block at $p$ is
  exactly one.  Lean checks this uniformly in
  \pdecl{factorialBlockPrivateQuotient_eq_one_of_gap_lt} and checks the
  two-owner consequence in
  \pdecl{factorialBlockFixedPairPrivateQuotients_eq_one}.  Thus the large
  private quotients seen at small blocks, for example the factor $719$ owned
  at $n=6$, are finite-range phenomena.  The factor-level reduction does not
  require two moving denominator indices: two factors inside one moving
  private quotient can suffice.  It still requires selected nontrivial factors
  that escape with $p$.
\end{itemize}

\begin{proposition}[fixed-owner absorption boundary]
\label{bdry:fixed-owner-absorption}
Fix denominator indices $i,j\geq2$.  Whenever
\[
 p>\max\{i!-1,j!-1\},
\]
both corresponding private quotients in the factorial block at $p$ equal
one.  Hence any cofinal private-factor certificate must use owner indices or
prime factors that escape with the block parameter.
\end{proposition}

\section{Finite certificates}\label{sec:finite}

The following are computations.  Each excludes exactly the denominators it
names and nothing more.

The finite-support vector $\lambda=2e_3-e_4$ has, by kernel check, $V_2=0$,
factorial moment $-12$, $V_3=-2$, $V_4=11$, and $V_d=-12$ for every $d\ge5$.
Under the exact rational tail enclosure $1/119<\Theta_4<1/50$, its residual lies
strictly between $-93/575$ and $-309/13685$; in particular it is nonzero and
subunit.

The combination of support $n\ge2$ with moment exactly $\Lc_D$ is impossible:
the left-hand side is even and $\Lc_D$ is odd.  Stronger, Lean checks
$12\Lc_D\mid M$ on that support after channels $2$ through $D$ vanish.  Auxiliary
kernels with an index-$1$ coordinate can realise moment $\Lc_D$, but that is a
different domain; the unavailable $D\le12$ coefficient arrays are not
reconstructed here.  The fully specified vector $\lambda=2e_3-e_4$ is retained
above.  At $D=3$ the vector $c=(-40,55,-10,1)$ on the
support $(3,4,5,6)$ annihilates channels $2$ and $3$, has moment $600$, and
satisfies
\[
  0.09925341997208298<\Lc_3(c)<0.09925341997208300 .
\]
This excludes denominators dividing $600$.

There are two different computations through the same range.  An external
interval computation reports the stronger geometric statement that no
zero-branch event occurs at any $m\le100000$; its cited executable and source
digest were not supplied, so that zero-branch classification remains external
finite evidence.  The strict-successor carry computation used above is local
and independently regenerated: its exact source, GMP backend, canonical
data, and receipt digests are recorded in the research packet.  The two
claims must not be conflated.  The local certificate establishes
$b_{300000}\ne1$, and the checked theorem converts precisely that fact into
$q\nmid299999!$ (and hence $q\ge300000$).  The bracket-certified continued-fraction prefix independently
raises the finite size exclusion to $q\ge2^{39990}>10^{12038}$; its script and
receipt are named above, and it likewise does not change a quantifier.

All of these conclusions are finite.

\section{Complements and further questions}\label{sec:open}

Four questions remain.

On the channel side, strict nonvanishing by itself is only a denominator
screen.  The
uniform remote-support theorem of \S\ref{sec:translator} already produces the
low-channel kernel, a nonzero moment, and a residual in $[-1/2,1/2]$ beyond
every cutoff.  For a hypothetical reduced denominator $q$, a nonzero residual
only says that $q$ does not divide the associated determinant; the bound
$1/2$ does not put it below the rational lattice spacing $1/q$.  The scalar
question is therefore the following: for a
hypothetical reduced denominator $b$, for some $D$ and some nonzero $m$
divisible by $\lcm(b,\Lc_D)$ with $mT_D$ nonintegral, equivalently for $k$ with
$0<|mT_D+k|<1$.  The formerly proposed determinantal inequality
$0<|N(I,D)|<\Delta(I,D)$ is impossible: Laplace expansion along the appended
row gives $\Delta(I,D)\mid N(I,D)$, so nonzero $N$ has
$|N|\ge\Delta$.  This divisibility rules out the proposed inequality for every
choice of $I$ and $D$.

On the projection side, endpoint transport to the literal factorial block is
checked.  Cofinal prefix-private factor supply is also available: Wilson gives
an unconditional existence theorem, while the source-backed large-prime
theorem and least-hit bridge give the stronger canonical selection.  What is
not known is a cofinal family for which the unit pair $(1,q)$ satisfies both
the global complementary-residue bound and the local collision-core bound.
The factor cannot stabilise at fixed low-index data: the absorption theorem
above makes every such private quotient equal to one for large $p$.
Reflection further shows that a linear-size factor at an odd later index may
be a repeated hit in the collision core rather than private support.

On the prime-pole side, the exact local target is now
$R_{q,e}(M)\ne0\pmod q$ for a suitably controlled cofinal family.  Mere
existence of a large divisor or a bounded number of maximal hits is
insufficient, as the two exact three-hit cancellations show; their reciprocal
equalities are Lean checked.  Even
nonvanishing must then be coupled to the Archimedean strict-successor window;
it cannot by itself change the global rationality quantifier.

On the grid side, the exact target is to prove that
$m\nmid Z_m$ for arbitrarily large $m$ (equivalently $b_m\ne1$ cofinally).
The adjacent-unit window of \S\ref{sec:adjacent-unit-no-go} becomes universal
after the pair assumption; it does not by itself close every predecessor
restriction.
Prime-index subtargets remain useful:
prove that infinitely many primes $p$ fail
$p\mid N_p$, or that for some fixed $k$ infinitely many primes $p$ fail
$p^k\mid N_{kp}$, where $N_m$ is the least integer strictly above $m!$ times the
$m$-th partial sum; prove for cofinally many primes $p$ that the least
$((p-1)!)^{-1}$ grid point strictly above $H_p$ is not $\Ser$; and beyond the
finite range already discharged by $q\nmid299999!$, control cancellation in
the first-crossing gap $V_p$ strongly enough to contradict
$\operatorname{den}(V_p)\ge\tau_p!-1$.  The prime-grid statement through
$p\le300000$ is a corollary of the finite carry certificate.

Two further questions are structural rather than immediate.  Find a global
invariant of the true divisor-factorial coefficients that rules out eventual
$q_m=1$ and is absent from the synthetic congruence countermodels; the
countermodels of \S\ref{sec:nogo} exist precisely because no such invariant is
yet in hand.  And construct or rule out an $N$-dependent, denominator-resonant,
strict-subrank vector Hermite--Pad\'e system that does not factor through scalar
evaluation at $z=1$ and does not reassemble the full factorial-gap lcm.  The
available calculations rule out the scalar and rank-saturated cases.  The
strict-subrank vector-valued case remains open.

Erd\H{o}s \#68 remains open: no statement above proves irrationality or
excludes every rational value.  The finite checked consequences are nevertheless
unconditional: every rational representation has
$q\ge2^{39990}>10^{12038}$ by the bracket-certified continued-fraction
screen; the independent carry certificate separately supplies
$q\nmid299999!$.

\section*{Statements and declarations}

The Lean declarations cited below are the proof authority for the formalized
statements.  The checked core comprises the canonical factorial digit kernel,
the finite defect automaton
algebra, floor-factorial channel arithmetic, the channel congruence and its
integral normal form, the prime unit translator, weighted projection rigidity,
the factor-split projection reduction, the fixed-index factorial-base
absorption no-go,
the rational-grid plateau and first-exit results, the first-crossing denominator
bounds, and the literal-prefix prime obstruction through the exact $p=11$
instance, strengthened by the all-index eventual-unit-carry theorem and the
exact reductions at $m=60,64,67$, the bound $q\ge67$, and the finite geometric
peeling identity.  It also checks the adjacent-unit width-one criterion, the
two-step denominator telescope and offset factorisation, and their reduction
to the universal numerator bound.  It checks the normalized strict-successor step and its
finite factorial-series expansion in the carry defects $1-b_m$, the convergence
$Z_m/m!\to\Ser$, and the exact equivalence between irrationality and cofinally
many non-unit carries.  The source-current Comparator front-loads the fixed
companion orbit as \texttt{companionOrbit\_completeCharacterization}, which
repeats Theorem~\ref{res:companion-orbit-rationality-boundary} and
Equation~\eqref{eq:companion-cofinal} literally.  Its second theorem packages
the complete four-way strict-successor result, namely eventual unit carry, cofinal
non-unit carry, pointwise carry/divisibility equivalence, and cofinal
strict-successor divisibility failure, as
\texttt{strictSuccessorCarry\_completeCharacterization}, routed to
Theorem~\ref{res:strict-successor-complete-characterization}.  The finite second-layer decomposition, its digit
identity $d^{(2)}_m=m-1-b_m$, the two-scale floor-stability theorem, and the
eventually-maximal rational-digit exclusion are checked separately in
\path{SecondLayerDigit.lean}.  The maximal prime-pole numerator formula and
the explicit reciprocal equalities at $139$ and $2593$ are checked in
\path{PrimePoleCriterion.lean}; the exact scan supplies the maximal-hit data
and resulting full cancellations.  These two source-current modules sharpen
the finite calculations, but neither proves a cofinal condition.  The exact
GMP carry certificate through $m=300000$ is separately regenerated and
hash-bound; combined with the checked carry theorem it gives
$q\nmid299999!$ (and hence $q\ge300000$),
but it is not itself a Lean evaluation.  The certified continued-fraction
script and receipt provide the independent finite bound
$q\ge2^{39990}>10^{12038}$.  This exact bracket computation is independent of
Lean and is finite.  The digit--rationality equivalence, the weighted primitive support
decomposition, and the determinant-quotient reduction are returned derivations
that have not been kernel-checked here, and are labelled as such wherever they
appear.  The factorial-gap lcm growth bound is derived here from
the exact factorial-congruence multiplicity theorem of Garaev--Luca--Shparlinski
\cite[arXiv v1, Thm.~12, p.~16]{garaev-luca-shparlinski}.  We derive the bound
from their theorem; it has not been formalized in Lean and no preceding result
depends on it.

\appendix
\section{Guide to the formal sources}\label{app:sources}

\begingroup
\raggedright
The public \texttt{ErdosProblems.Erdos68} package contains the checked source
for this note.  The pinned release snapshot used by the declaration links
contains the following core modules:
\texttt{CanonicalFactorialDigits},
\texttt{ChannelBreakpointRigidity},
\texttt{ChannelIntegralCongruence},
\texttt{CompanionOrbitRationality},
\texttt{EndpointWeightedPrivateSupport},
\texttt{FactorialCarry},
\texttt{FactorialChannelCertificate},
\texttt{FactorialZeroPlateau},
\texttt{FiniteDefectAutomaton},
\texttt{PrimeUnitTranslator},
\texttt{PrimeZeroBranch}, and
\texttt{StrictSuccessorArithmetic}.
The release root imports every cited module.  The declaration table below is
pinned to the shared formal-source commit used throughout this problem-note
series.\par
\endgroup

\begin{itemize}[leftmargin=*]
\item \lref{Erdos68/FactorialDigitRigidity.lean}{21}{strictSuccessor}
\item \lref{Erdos68/FactorialDigitRigidity.lean}{27}{strictSuccessor_sub_int}
\item \lref{Erdos68/FactorialDigitRigidity.lean}{34}{canonicalFactorialDigit}
\item \lref{Erdos68/FactorialDigitRigidity.lean}{40}{factorial_scaled_rational_eq_intCast}
\item \lref{Erdos68/FactorialDigitRigidity.lean}{63}{canonicalFactorialDigit_eq_zero_of_rational}
\item \lref{Erdos68/FactorialDigitRigidity.lean}{92}{prime_pow_dvd_factorial_dilation}
\item \lref{Erdos68/FactorialDigitRigidity.lean}{102}{prime_pow_dvd_factorial_dilation_div}
\item \lref{Erdos68/FactorialDigitRigidity.lean}{115}{prime_pow_dvd_cleared_rational_numerator}
\item \lref{Erdos68/DivisorFactorialCentre.lean}{17}{residualCentreTerm}
\item \lref{Erdos68/DivisorFactorialCentre.lean}{22}{factorialCoeffTerm}
\item \lref{Erdos68/DivisorFactorialCentre.lean}{27}{residualCentre}
\item \lref{Erdos68/DivisorFactorialCentre.lean}{32}{factorialCoeffRat}
\item \lref{Erdos68/DivisorFactorialCentre.lean}{35}{factorial_eq_mul_pred_factorial}
\item \lref{Erdos68/DivisorFactorialCentre.lean}{43}{pred_div_eq_of_not_dvd}
\item \lref{Erdos68/DivisorFactorialCentre.lean}{57}{pred_div_add_one_eq_of_dvd}
\item \lref{Erdos68/DivisorFactorialCentre.lean}{84}{residualCentreTerm_of_not_dvd}
\item \lref{Erdos68/DivisorFactorialCentre.lean}{101}{residualCentreTerm_of_dvd}
\item \lref{Erdos68/DivisorFactorialCentre.lean}{125}{residualCentreTerm_self}
\item \lref{Erdos68/DivisorFactorialCentre.lean}{134}{factorialCoeffTerm_self}
\item \lref{Erdos68/DivisorFactorialCentre.lean}{143}{residualCentre_recurrence}
\item \lref{Erdos68/FactorialChannelCertificate.lean}{15}{factorial_pow_floor_dvd_factorial}
\item \lref{Erdos68/FactorialChannelCertificate.lean}{29}{channelWeight}
\item \lref{Erdos68/FactorialChannelCertificate.lean}{33}{channelWeight_mul_denominator}
\item \lref{Erdos68/FactorialChannelCertificate.lean}{41}{channelNumerator}
\item \lref{Erdos68/FactorialChannelCertificate.lean}{45}{factorialMoment}
\item \lref{Erdos68/FactorialChannelCertificate.lean}{49}{channelEvent}
\item \lref{Erdos68/FactorialChannelCertificate.lean}{55}{channelEvent_eq_zero_of_not_dvd}
\item \lref{Erdos68/FactorialChannelCertificate.lean}{91}{lambda34}
\item \lref{Erdos68/FactorialChannelCertificate.lean}{95}{lambda34_apply_three}
\item \lref{Erdos68/FactorialChannelCertificate.lean}{99}{lambda34_apply_four}
\item \lref{Erdos68/FactorialChannelCertificate.lean}{102}{lambda34_channel_two}
\item \lref{Erdos68/FactorialChannelCertificate.lean}{109}{lambda34_factorialMoment}
\item \lref{Erdos68/FactorialChannelCertificate.lean}{116}{lambda34_channel_three}
\item \lref{Erdos68/FactorialChannelCertificate.lean}{123}{lambda34_channel_four}
\item \lref{Erdos68/FactorialChannelCertificate.lean}{130}{lambda34_channel_ge_five}
\item \lref{Erdos68/FactorialChannelCertificate.lean}{141}{lambda34ResidualOfTail}
\item \lref{Erdos68/FactorialChannelCertificate.lean}{146}{lambda34_residual_bounds}
\end{itemize}

\begingroup
\raggedright
\paragraph{Source-current adjacent-unit no-go.}
Section~\ref{sec:adjacent-unit-no-go} is checked in the local module
\path{ErdosProblems/Erdos68/AdjacentUnitCarryWindow.lean}, principally through
\path{consecutive_unit_carries_iff_positive_offset_le_den},
\path{twoStep_den_mul_transitionNormalizers},
\path{adjacentUnitCarryWindowDen_eq_twoStep_den}, and
\path{adjacentUnitCarryWindowOffset_eq_twoStep_factorization}.  The universal numerator
bounds are in \path{PrimeZeroBranch.lean}.  This module lies outside the older
public declaration-table pin and must accompany the revised paper at its next
public checkpoint.
\par
\medskip
\paragraph{Source-current second-layer and prime-pole modules.}
The finite digit law and two-scale floor-stability boundary are checked in
\path{SecondLayerDigit.lean}; the general prime-pole formula and the $139$ and
$2593$ reciprocal equalities are checked in \path{PrimePoleCriterion.lean}.
Both are at source commit
\texttt{25e960b69ea1346ffc511b199abf4cc755eb6df8}.  The local exact scan supplies
the complete maximal-hit payloads and resulting cancellations.  Source paths
and result identifiers are in \path{research_packet.json}.  These local
identities and the manuscript revision require the same public checkpoint.
\par
\medskip
\paragraph{Source-current fixed companion orbit.}
The complete infinite rationality boundary is checked in
\path{CompanionOrbitRationality.lean}.  Its paper-facing endpoints are
\path{not_irrational_factorialGapSeries_iff_eventually_companion_floor_neg_two}
and
\path{irrational_factorialGapSeries_iff_cofinal_companion_floor_misses}; the
generic fixed-orbit theorem, canonical factorial expansion, positive sub-unit
exponential tail, and reverse telescoping assembly are in the same module.
The focused source build and axiom print are green.  This source-current module,
the two-theorem Comparator revision, and this manuscript revision must be
published at one common immutable checkpoint before external replay or Palomar
submission readiness is claimed.
\par
\endgroup

\clearpage
\begin{thebibliography}{9}
\small
\raggedright
\setlength{\itemsep}{0pt}
\bibitem{erdos1988} P.~Erd\H{o}s, \emph{On the irrationality of certain series:
  problems and results}, in A.~Baker (ed.), \emph{New Advances in Transcendence
  Theory}, Cambridge UP, 1988, pp.~102--109,
  doi:\href{https://doi.org/10.1017/CBO9780511897184.009}
  {10.1017/CBO9780511897184.009}.
\bibitem{bloom} T.~F. Bloom, \emph{Erd\H{o}s Problem \#68}.
  \url{https://www.erdosproblems.com/68}, accessed 28 July 2026.
\bibitem{garaev-luca-shparlinski}
  M.~Z. Garaev, F.~Luca, and I.~E. Shparlinski,
  \emph{Character sums and congruences with $n!$},
  Trans.\ Amer.\ Math.\ Soc.\ \textbf{356} (2004), no.~12, 5089--5102.
  \url{https://doi.org/10.1090/S0002-9947-04-03612-8};
  arXiv:\href{https://arxiv.org/abs/math/0403422}{math/0403422v1}.
\bibitem{stewart2004}
  C.~L. Stewart,
  \emph{On the greatest and least prime factors of $n!+1$, II},
  Publ.\ Math.\ Debrecen \textbf{65} (2004), no.~3--4, 461--480.
  \url{https://publi.math.unideb.hu/paper/989/download/10_5486_PMD_2004_3190.pdf}.
\bibitem{koepf-schmersau}
  W.~Koepf and D.~Schmersau,
  \emph{Irrationality of certain infinite series II},
  Analysis \textbf{31} (2011), 117--124.
  \url{https://doi.org/10.1524/anly.2011.1094}.
\bibitem{duverney}
  D.~Duverney,
  \emph{Irrationality of fast converging series of rational numbers},
  J.\ Math.\ Sci.\ Univ.\ Tokyo \textbf{8} (2001), 275--316.
  \url{https://www.ms.u-tokyo.ac.jp/journal/pdf/jms080206.pdf}.
\bibitem{cantor1869} G.~Cantor, \emph{\"Uber die einfachen Zahlensysteme},
  Z. Math. Phys. \textbf{14} (1869), 121--128.
\bibitem{galambos1976} J.~Galambos, \emph{Representations of Real Numbers by
  Infinite Series}, Lecture Notes in Math. \textbf{502}, Springer, 1976,
  Chapter~1.
\bibitem{hancl-tijdeman}
  J.~Han\v{c}l and R.~Tijdeman,
  \emph{On the irrationality of factorial series},
  Acta Arith.\ \textbf{118} (2005), no.~4, 383--401.
  \url{https://www.impan.pl/shop/en/publication/transaction/download/product/83588}.
\bibitem{barreto-et-al}
  K.~Barreto, J.~Kang, S.-h.~Kim, V.~Kova\v{c}, and S.~Zhang,
  \emph{Irrationality of rapidly converging series: a problem of Erd\H{o}s
  and Graham},
  arXiv:\href{https://arxiv.org/abs/2601.21442v3}{2601.21442v3}, 2026.
\bibitem{louwsma-martino}
  J.~Louwsma and J.~Martino,
  \emph{Rational numbers with odd greedy expansion of fixed length},
  arXiv:\href{https://arxiv.org/abs/2309.07280}{2309.07280v1}, 2023,
  Lemma~4.1, p.~10.
\bibitem{lai}
  L.~Lai, \emph{On the largest prime divisor of $n!+1$},
  arXiv:\href{https://arxiv.org/abs/2103.14894}{2103.14894v1}, 2021,
  proof of Lemma~2.4, (2.5).
\end{thebibliography}

\end{document}
